\section*{Bab \ref{ch:g6:matter-from-food} --- Materi dari Makanan}

\begin{solution}{exo:g6:matter-from-food:1}
\omterm{def:g5:food-webs:roles}{Konsumen} mengubah \omterm{def:g6:plants-are-producers:matter}{materi organik} yang dimakannya menjadi \omterm{def:g6:plants-are-producers:matter}{materi
organiknya} sendiri. \omterm{def:g5:food-webs:roles}{Produsen} membuat \omterm{def:g6:plants-are-producers:matter}{materi organik} yang baru dari
pasokan yang murni mineral, dengan \omterm{prop:g3:balanced-diet:jobs}{energi} \omterm{def:g6:exploring-our-environment:conditions}{cahaya}.
\end{solution}

\begin{solution}{exo:g6:matter-from-food:2}
\omterm{def:g4:journey-of-food:digestion}{Pencernaan} membongkar rumputnya menjadi \omterm{def:g4:journey-of-food:digestion}{zat gizi}
(\cref{ch:g4:journey-of-food}); \omterm{prop:g4:journey-of-food:absorb}{darahnya} mengantarkannya ke mana-mana
(\cref{ch:g4:heart-and-blood}); \omterm{def:g5:first-look-at-cells:cell}{sel} tubuhnya merangkainya kembali
menjadi materi-sapi, sambil membelah seiring tumbuhnya tubuhnya
(\cref{ch:g5:first-look-at-cells}).
\end{solution}

\begin{solution}{exo:g6:matter-from-food:3}
Pertumbuhan (kilogram tambahan anak sapinya); pembaruan (\omterm{def:g1:animals-around-us:covering}{bulu} yang
ditanggalkan, \omterm{def:g1:animals-around-us:covering}{bulu burung} baru, ranggah); perbaikan (\omterm{def:g3:bones-and-muscles:skeleton}{tulang} yang
menyatu); hasil (\omterm{def:g2:animals-grow-up:born}{susu}, \omterm{def:g2:animals-grow-up:born}{telur}, wol).
\end{solution}

\begin{solution}{exo:g6:matter-from-food:4}
Tidak pernah terserap --- keluar sebagai tinja; dibakar dengan oksigen
dari napasnya untuk gerakan dan kehangatan --- keluar sebagai gas yang
diembuskan; dan sisanya, yang dibangun menjadi pertumbuhan dan hasil.
\end{solution}

\begin{solution}{exo:g6:matter-from-food:5}
Karena \qty{25}{kg} \omterm{def:g2:animals-grow-up:born}{susu} tiap harinya adalah \omterm{def:g6:plants-are-producers:matter}{materi organik} yang harus
dibangunnya ulang dari pakan hari itu; \omterm{def:g2:animals-grow-up:born}{susu} adalah produksi, dan
produksi adalah makanan yang dibangun ulang --- tanpa masukan, tanpa
keluaran.
\end{solution}

\begin{solution}{exo:g6:matter-from-food:6}
Dirangkai \emph{di dalam} ayamnya, dari \omterm{def:g4:journey-of-food:digestion}{zat gizi} biji-bijiannya ---
tetapi \omterm{def:g6:plants-are-producers:matter}{materi organiknya} sendiri pertama kali dibuat sebelum ayamnya,
di dalam \omterm{def:g1:plants-around-us:parts}{daun} tanaman biji-bijiannya. Ayamnya mengubah; \omterm{def:g1:plants-around-us:plant}{tumbuhannya}
menghasilkan.
\end{solution}

\begin{solution}{exo:g6:matter-from-food:7}
Pada tiap anak panahnya, sebagian besar materi yang dimakan dibakar
untuk hidup atau hilang sebagai tinja, dan hanya pecahan yang sederhana
yang menjadi tubuh si pemakan --- satu-satunya bagian yang dapat
dimakan lantai berikutnya. Karena itu tiap lantainya hanya punya
sebuah pecahan dari materi di bawahnya: piramidanya menyempit oleh
hitung-hitungan.
\end{solution}

\begin{solution}{exo:g6:matter-from-food:8}
Masuk: daging atau makanan kering (dan sesekali seekor tikus).
Dibangun: pertumbuhannya sebagai anak kucing, \omterm{def:g1:animals-around-us:covering}{bulu} yang diperbaruinya
pada tiap pergantian \omterm{def:g1:animals-around-us:covering}{bulu}, perbaikan sesudah lecet. Keluar: tinja di
kebun, dan bagian yang terbakar yang menenagai setiap terkaman,
dengkuran dan \omterm{prop:g1:healthy-habits:needs}{tidur} siang yang hangat.
\end{solution}

\begin{solution}{exo:g6:matter-from-food:9}
Yang berbeda adalah sumbernya dan perlengkapannya: rumput lawan
kelinci, penggilas-dan-usus-luas lawan penerkam-dan-usus-pendek. Yang
serupa adalah perdagangannya sendiri: keduanya membongkar \omterm{def:g6:plants-are-producers:matter}{materi
organik} yang dimakan menjadi \omterm{def:g4:journey-of-food:digestion}{zat gizi} lalu membangunnya ulang sebagai
zat tubuhnya sendiri.
\end{solution}

\begin{solution}{exo:g6:matter-from-food:10}
Sebagian besarnya keluar sebagai bagian yang terbakar --- yang
menenagai setahun berlari, berpikir dan bertahan pada
\qty{37}{\degreeCelsius} --- dan sebagian tidak pernah terserap lalu
keluar sebagai tinja. Hanya sisanya yang kecil, yaitu
\qty{3}{kg}-nya, yang menjadi tubuh yang baru: pembukuan anaknya
seimbang seperti pembukuan sapinya.
\end{solution}

\begin{solution}{exo:g6:matter-from-food:11}
Tubuh siputnya menyaring bahan mineral \omterm{def:g1:animals-around-us:covering}{cangkangnya} dari makanannya dan
sekelilingnya, mengangkutnya, lalu mengendapkannya lapis demi lapis
dalam arsitekturnya sendiri. Zatnya mineral; \emph{pembangunannya}
adalah pekerjaan \omterm{def:g1:animals-around-us:animal}{hewan} yang hidup --- sebuah benda buatan, dari bahan
mineral.
\end{solution}

\begin{solution}{exo:g6:matter-from-food:12}
Kalau dijalurkan lewat babinya, sebagian besar materi kentangnya
dibelanjakan untuk kehidupan babinya sendiri --- dibakar untuk
kehangatan dan gerakan, hilang sebagai tinja --- dan hanya sebuah
pecahan yang kembali sebagai daging babi. Kalau dimakan langsung,
kentangnya melewatkan pengeluaran babinya. Perbedaannya naik ke
pembukuan babinya, bukan ke atas mejanya.
\end{solution}

\begin{solution}{pb:g6:matter-from-food:1}
\textbf{1.} Rerumputan padang penggembalaannya dan tanaman petak
sayurnya (serta pohon \omterm{prop:g3:flowers-fruits-seeds:transform}{buah} yang ada). Masukannya: air dan \omterm{prop:g4:what-plants-need:needs}{garam
mineral} dari tanah, karbon dioksida dari udara. \omterm{prop:g3:balanced-diet:jobs}{Energinya}: \omterm{def:g6:exploring-our-environment:conditions}{cahaya}.

\textbf{2.} Di dalam \omterm{def:g1:plants-around-us:parts}{daun} tanaman rumputnya sendiri --- setiap
gramnya: \omterm{def:g5:food-webs:roles}{produsen} membuat di tempat, pada siang hari.

\textbf{3.} Kembalikan apa yang kamu pinjam: pelapukan pupuk kandangnya
mengisi ulang \omterm{prop:g4:what-plants-need:needs}{garam mineral} tanah petaknya, dan menggantikan apa yang
diangkut pergi oleh panen.

\textbf{4.} Karena sapinya adalah seorang \omterm{def:g5:food-webs:roles}{konsumen}: perdagangannya
membutuhkan materi \emph{organik} untuk dibongkar lalu dibangun ulang.
\omterm{prop:g4:what-plants-need:needs}{Garam mineral} dan air adalah masukan seorang \omterm{def:g5:food-webs:roles}{produsen}; ia tidak punya
permesinan untuk membangun tubuh dari mineral --- justru perdagangan
itulah yang tidak dimilikinya.

\textbf{5.} Sebagian melintas tanpa terserap --- kotoran sapinya;
sebagian besar dibakar dengan oksigen yang dihirup untuk menenagai
berjalan, mengunyah dan bertahan hangat; sisanya menjadi \omterm{def:g2:animals-grow-up:born}{susu} dan sapi
(dan pembukuan hari itu pun seimbang).

\textbf{6.} Masuk: biji-bijian. Dibangun: \omterm{def:g2:animals-grow-up:born}{telur}, pemeliharaan dirinya
sendiri, dan --- ketika berganti \omterm{def:g1:animals-around-us:covering}{bulu} --- satu mantel \omterm{def:g1:animals-around-us:covering}{bulu burung} yang
sama sekali baru. Keluar: tinja dan bagian yang terbakar. Ketika
berganti \omterm{def:g1:animals-around-us:covering}{bulu}, baris ``dibangun''-nya dibelanjakan untuk \omterm{def:g1:animals-around-us:covering}{bulunya};
dengan masukan yang tidak berubah, catatan \omterm{def:g2:animals-grow-up:born}{telurnya} harus menyusut
untuk sementara waktu.

\textbf{7.} Rumput $\rightarrow$ ayam (lewat biji-bijian, yang sendiri
merupakan hasil \omterm{def:g1:plants-around-us:plant}{tumbuhan}) $\rightarrow$ anak ayam $\rightarrow$ elang;
atau di sisi padang penggembalaannya: rumput $\rightarrow$ sapi.
Materinya \emph{dibuat} hanya pada mata rantai hijaunya (rumput,
tanaman biji-bijian); setiap mata rantai \omterm{def:g1:animals-around-us:animal}{hewan} --- ayam, anak ayam,
elang, sapi --- hanya mengubahnya.

\textbf{8.} Roti: \omterm{def:g1:plants-around-us:parts}{daun} tanaman gandumnya. \omterm{def:g2:animals-grow-up:born}{Telur} rebus: ayam dari
biji-bijian --- \omterm{def:g1:plants-around-us:parts}{daun} tanaman biji-bijiannya. \omterm{def:g2:animals-grow-up:born}{Susu}: sapi dari rumput
--- \omterm{def:g1:plants-around-us:parts}{daun} rumputnya. Tiga jejak, tiga \omterm{def:g1:plants-around-us:parts}{daun}, semuanya di dalam \omterm{def:g6:exploring-our-environment:conditions}{cahaya}.

\textbf{9.} Tiap lantai hanya meneruskan sebuah pecahan: dari
\qty{10000}{kg} rumputnya, sebuah populasi kelinci hanya membangun
pecahannya yang sederhana sebagai tubuh; dari jumlah itu, rubah hanya
membangun sebuah pecahan lagi. Dua lantai kehilangan menyusutkan
rumput senilai satu padang penggembalaan menjadi rubah senilai satu
sarang --- perkiraannya adalah hitung-hitungan piramidanya.

\textbf{10.} Kalau dimakan langsung, hasil petaknya sampai kepada
keluarganya tanpa satu lantai \omterm{def:g5:food-webs:roles}{konsumen} pun di antaranya: tidak ada
bagiannya yang dibakar dalam kehidupan seekor \omterm{def:g1:animals-around-us:animal}{hewan} atau hilang dalam
tinjanya. Jalur padang penggembalaannya membayar seluruh pengeluaran
sapinya lebih dahulu, sehingga tiap meter perseginya memberi makan
lebih sedikit orang.

\textbf{11.} Rumputnya berhenti menghasilkan, sehingga masukan sapinya
gagal: mula-mula \omterm{def:g2:animals-grow-up:born}{susunya} menyusut (produksi adalah makanan yang
dibangun ulang), lalu ia membakar cadangan lemaknya sendiri dan
kehilangan berat --- baris ``dibangun''-nya menjadi negatif. Sinar
matahari adalah \omterm{prop:g3:balanced-diet:jobs}{energi}, bukan zat: tanpa air tidak ada yang dapat
dibangun rumputnya, sehingga padang penggembalaannya tidak dapat
diselamatkan dari langit.

\textbf{12.} \omterm{def:g5:food-webs:roles}{Produsen}: \omterm{def:g1:plants-around-us:plant}{tumbuhan} di lahannya membuat seluruh \omterm{def:g6:plants-are-producers:matter}{materi
organiknya} dari pasokan mineral dan \omterm{def:g6:exploring-our-environment:conditions}{cahaya}. \omterm{def:g5:food-webs:roles}{Konsumen}: sapi, ayam dan
keluarganya mengubah materi itu, membakar sebagian besarnya lalu
membangun sisanya. Yang hilang: para \omterm{def:g5:food-webs:roles}{pengurainya} --- yang harus
mengembalikan tinjanya, kotoran sapinya dan sisa musim gugurnya
menjadi \omterm{def:g6:plants-are-producers:matter}{materi mineral}; merekalah pokok bab berikutnya.
\end{solution}
